% Part: normal-modal-logic
% Chapter: tableaux
% Section: countermodels

\documentclass[../../../include/open-logic-section]{subfiles}

\begin{document}

\olfileid{nml}{tab}{cou}
      
\olsection{Kontramodel saka \usetoken{P}{tableau}}

\begin{explain}
  Bukti teorema kelengkapan ora mung nuduhake yen saka
  $\Entails !A$ bisa dipikolehi $\Proves !A$, nanging uga maringi
  cara kanggo mbangun kontramodel tumrap~$!A$ yen $\Entails/ !A$.
  Kanggo~$\Log{K}$, cara iki minangka \emph{prosedur kaputusan}.
  Upamane $\Entails/ !A$. Bukti
  \olref[cpl]{prop:complete-tableau} maringi cara mbangun
  !!{tableau} jangkep. Cara iki tansah rampung. Aturan proposisional
  kanggo~$\Log{K}$ mung nambahake !!{formula}s mawa prefiks kang
  kompleksitase luwih cendhek; tegese saben aturan proposisional
  cukup ditrapake sapisan ing sawijine cabang tumrap saben rumus
  mawa tandha $\sFmla{S}{!A}[\sigma]$. Prefiks anyar mung diasilake
  dening aturan
  \iftag{prvBox}{$\TRule{\False}{\Box}$}{}\iftag{notprvBox,notprvDiamond}{}{
    and }\iftag{prvDiamond}{$\TRule{\True}{\Diamond}$}{}
  \iftag{notprvBox,notprvDiamond}{iki}{iki}, lan saben aturan mau uga
  cukup ditrapake sapisan (ngasilake siji prefiks anyar).
  \iftag{prvBox}{$\TRule{\True}{\Box}$}{}\iftag{notprvBox,notprvDiamond}{}{
    and }\iftag{prvDiamond}{$\TRule{\False}{\Diamond}$}{}
  \iftag{notprvBox,notprvDiamond}{kadhang}{kadhang} kudu ditrapake
  kaping pirang-pirang, nanging mung sapisan kanggo saben prefiks;
  prefiks anyar kang diasilake uga winates cacahe. Mula, sawisé
  tataran kang winates, pambangunan iki ngasilake cabang katutup
  utawa cabang jangkep.

  Yen wis kabangun tableau kang nduweni cabang bukak lan jangkep,
  bukti \olref[cpl]{thm:tableau-completeness} maringi model eksplisit
  kang nyukupi himpunan !!{formula}s mawa prefiks wiwitan. Dadi,
  kajaba kasunyatan yen $\Gamma \Entails !A$ njalari anane
  !!{tableau} katutup lan $\Gamma \Proves !A$, nalika kita nggoleki
  !!{tableau} katutup kanthi cara kang trep nanging nemokake cabang
  bukak lan jangkep ing !!{tableau}, kita ora mung ngerti yen $\Gamma \Entails/ !A$,
  nanging uga bisa mbangun kontramodel.
\end{explain}

\iftag{prvBox}{
\begin{ex}
  Kita ngerti yen $\Proves/ \Box(p \lor q) \lif (\Box p \lor \Box
  q)$. Pambangunan tableau diwiwiti mangkene:
  \begin{oltableau}
    [\pFmla{\False}{\Box(p \lor q) \lif (\Box p \lor \Box q)}{1},
      just = \TAss, checked
      [\pFmla{\True}{\Box(p \lor q)}{1},
        just = {\TRule{\False}{\lif}[1]}, 
        [\pFmla{\False}{\Box p \lor \Box q}{1},
          just = {\TRule{\False}{\lif}[1]}, checked
          [\pFmla{\False}{\Box p}{1},
            just = {\TRule{\False}{\lor}[3]}, checked
            [\pFmla{\False}{\Box q}{1},
              just = {\TRule{\False}{\lor}[3]}, checked
              [\pFmla{\False}{p}{1.1}, 
                just = {\TRule{\False}{\Box}[4]}, checked
                [\pFmla{\False}{q}{1.2}, 
                  just = {\TRule{\False}{\Box}[5]}, checked
                ]
              ]
            ]
          ]
        ]
      ]
    ]
  \end{oltableau}
  !!{tableau} iki mesthi durung rampung. Ing langkah sabanjure,
  kita gatekake siji-sijine larik kang durung diwenehi tandha centhang:
  !!{formula} mawa prefiks $\sFmla{\True}{\Box(p \lor q)}[1]$
  ing larik~$2$. Kanggo mbangun tableau jangkep, aturan
  $\TRule{\True}{\Box}$ ditrapake kanggo saben prefiks kang wis
  dienggo ing cabang iki, yaiku $1.1$ lan $1.2$:
  \begin{oltableau}
    [\pFmla{\False}{\Box(p \lor q) \lif (\Box p \lor \Box q)}{1},
      just = \TAss, checked
      [\pFmla{\True}{\Box(p \lor q)}{1},
        just = {\TRule{\False}{\lif}[1]}, 
        [\pFmla{\False}{\Box p \lor \Box q}{1},
          just = {\TRule{\False}{\lif}[1]}, checked
          [\pFmla{\False}{\Box p}{1},
            just = {\TRule{\False}{\lor}[3]}, checked
            [\pFmla{\False}{\Box q}{1},
              just = {\TRule{\False}{\lor}[3]}, checked
              [\pFmla{\False}{p}{1.1}, 
                just = {\TRule{\False}{\Box}[4]}, checked
                [\pFmla{\False}{q}{1.2}, 
                  just = {\TRule{\False}{\Box}[5]}, checked
                  [\pFmla{\True}{p \lor q}{1.1}, 
                    just = {\TRule{\True}{\Box}[2]}
                    [\pFmla{\True}{p \lor q}{1.2}, 
                      just = {\TRule{\True}{\Box}[2]}
                    ]
                  ]
                ]
              ]
            ]
          ]
        ]
      ]
    ]
  \end{oltableau}
  Saiki larik 2, 8, lan 9 durung diwenehi tandha centhang. Nanging
  ora ana prefiks anyar, mula kita trapake $\TRule{\True}{\lor}$
  marang larik~8 lan~9 ing kabeh cabang asile kang durung katutup:
  \begin{oltableau}
    [\pFmla{\False}{\Box(p \lor q) \lif (\Box p \lor \Box q)}{1},
      just = \TAss, checked
      [\pFmla{\True}{\Box(p \lor q)}{1},
        just = {\TRule{\False}{\lif}[1]}, checked
        [\pFmla{\False}{\Box p \lor \Box q}{1},
          just = {\TRule{\False}{\lif}[1]}, checked
          [\pFmla{\False}{\Box p}{1},
            just = {\TRule{\False}{\lor}[3]}, checked
            [\pFmla{\False}{\Box q}{1},
              just = {\TRule{\False}{\lor}[3]}, checked
              [\pFmla{\False}{p}{1.1}, 
                just = {\TRule{\False}{\Box}[4]}, checked
                [\pFmla{\False}{q}{1.2}, 
                  just = {\TRule{\False}{\Box}[5]}, checked
                  [\pFmla{\True}{p \lor q}{1.1}, 
                    just = {\TRule{\True}{\Box}[2]}, checked
                    [\pFmla{\True}{p \lor q}{1.2}, 
                      just = {\TRule{\True}{\Box}[2]}, checked
                      [\pFmla{\True}{p}{1.1},
                        just = {\TRule{\True}{\lor}[8]}, checked, close
                      ]
                      [\pFmla{\True}{q}{1.1},
                        just = {\TRule{\True}{\lor}[8]}, checked
                        [\pFmla{\True}{p}{1.2},
                          just = {\TRule{\True}{\lor}[9]}, checked]
                        [\pFmla{\True}{q}{1.2},
                          just = {\TRule{\True}{\lor}[9]}, checked, close]
                      ]
                    ]
                  ]
                ]
              ]
            ]
          ]
        ]
      ]
    ]
  \end{oltableau}
  Kari siji cabang bukak, lan cabang iku jangkep. Saka cabang iki
  kita netepake model kanthi himpunan jagad $W = \{1, 1.1, 1.2\}$
  (mung prefiks iki kang katon ing cabang bukak), relasi
  aksesibilitas $R = \{\tuple{1, 1.1}, \tuple{1, 1.2}\}$, lan
  panetepan $V(p) = \{1.2\}$ (amarga larik~12 ngemot
  $\sFmla{\True}{p}[1.2]$) sarta $V(q) = \{1.1\}$ (amarga larik~11
  ngemot $\sFmla{\True}{q}[1.1]$). Model iki kagambar ing
  \olref{fig:counter-Box}; bisa dipriksa yen iki kontramodel tumrap
  $\Box(p \lor q) \lif (\Box p \lor \Box q)$.
  \begin{figure}
  \begin{center}
    \begin{tikzpicture}[modal]
      \node[world] (w1) [label={[align=right]right:\mFalse{p}\\ \mFalse{q}}]{$1$}; 
      \node[world] (w2) [label={[align=right]right:\mFalse{p}\\ \mTrue{q}},
        above left=of w1]{$1.1$}; 
      \node[world] (w3) [label={[align=right]right:\mTrue{p}\\ \mFalse{q}},
        above right=of w1] {$1.2$};
      \draw[->] (w1) to (w2);
      \draw[->] (w1) to (w3);
    \end{tikzpicture}
  \end{center}
  \caption{Kontramodel tumrap $\Box(p \lor q) \lif (\Box p \lor \Box
    q)$.}
  \ollabel{fig:counter-Box}
\end{figure}
\end{ex}          
}
{
\begin{ex}
  Kita ngerti yen $\Proves/ (\Diamond p \land \Diamond q) \lif \Diamond(p
  \land q)$. Pambangunan tableau diwiwiti mangkene:
  \begin{oltableau}
    [\pFmla{\False}{(\Diamond p \land \Diamond q) \lif \Diamond(p \land q)}{1},
      just = \TAss, checked
      [\pFmla{\True}{\Diamond p \land \Diamond q}{1},
        just = {\TRule{\False}{\lif}[1]}, checked
        [\pFmla{\False}{\Diamond(p \land q)}{1},
          just = {\TRule{\False}{\lif}[1]}
          [\pFmla{\True}{\Diamond p}{1},
            just = {\TRule{\True}{\land}[2]}, checked
            [\pFmla{\True}{\Diamond q}{1},
              just = {\TRule{\True}{\land}[2]}, checked
              [\pFmla{\True}{p}{1.1}, 
                just = {\TRule{\True}{\Diamond}[4]}, checked
                [\pFmla{\True}{q}{1.2}, 
                  just = {\TRule{\True}{\Diamond}[5]}, checked
                ]
              ]
            ]
          ]
        ]
      ]
    ]
  \end{oltableau}
  !!{tableau} iki mesthi durung rampung. Ing langkah sabanjure,
  kita gatekake siji-sijine larik kang durung diwenehi tandha centhang:
  !!{formula} mawa prefiks $\sFmla{\False}{\Diamond(p \land q)}[1]$
  ing larik~$3$. Kanggo mbangun tableau jangkep, aturan
  $\TRule{\False}{\Diamond}$ ditrapake kanggo saben prefiks kang wis
  dienggo ing cabang iki, yaiku $1.1$ lan $1.2$:
  \begin{oltableau}
    [\pFmla{\False}{(\Diamond p \land \Diamond q) \lif \Diamond(p \land q)}{1},
      just = \TAss, checked
      [\pFmla{\True}{\Diamond p \land \Diamond q}{1},
        just = {\TRule{\False}{\lif}[1]}, checked
        [\pFmla{\False}{\Diamond (p \land q)}{1},
          just = {\TRule{\False}{\lif}[1]}
          [\pFmla{\True}{\Diamond p}{1},
            just = {\TRule{\True}{\land}[2]}, checked
            [\pFmla{\True}{\Diamond q}{1},
              just = {\TRule{\True}{\land}[2]}, checked
              [\pFmla{\True}{p}{1.1}, 
                just = {\TRule{\True}{\Diamond}[4]}, checked
                [\pFmla{\True}{q}{1.2}, 
                  just = {\TRule{\True}{\Diamond}[5]}, checked
                  [\pFmla{\False}{p \land q}{1.1}, 
                    just = {\TRule{\False}{\Diamond}[3]}
                    [\pFmla{\False}{p \land q}{1.2}, 
                      just = {\TRule{\False}{\Diamond}[3]}
                    ]
                  ]
                ]
              ]
            ]
          ]
        ]
      ]
    ]
  \end{oltableau}
  Saiki larik 3, 8, lan 9 durung diwenehi tandha centhang. Nanging
  ora ana prefiks anyar, mula kita trapake $\TRule{\False}{\land}$
  marang larik~8 lan~9 ing kabeh cabang asile kang durung katutup:
  \begin{oltableau}
    [\pFmla{\False}{(\Diamond p \land \Diamond q) \lif \Diamond(p \land q)}{1},
      just = \TAss, checked
      [\pFmla{\True}{\Diamond p \land \Diamond q}{1},
        just = {\TRule{\False}{\lif}[1]}, checked
        [\pFmla{\False}{\Diamond(p \land q)}{1},
          just = {\TRule{\False}{\lif}[1]}, checked
          [\pFmla{\True}{\Diamond p}{1},
            just = {\TRule{\True}{\land}[2]}, checked
            [\pFmla{\True}{\Diamond q}{1},
              just = {\TRule{\True}{\land}[2]}, checked
              [\pFmla{\True}{p}{1.1}, 
                just = {\TRule{\True}{\Diamond}[4]}, checked
                [\pFmla{\True}{q}{1.2}, 
                  just = {\TRule{\True}{\Diamond}[5]}, checked
                  [\pFmla{\False}{p \land q}{1.1}, 
                    just = {\TRule{\False}{\Diamond}[3]}, checked
                    [\pFmla{\False}{p \land q}{1.2}, 
                      just = {\TRule{\False}{\Diamond}[3]}, checked
                      [\pFmla{\False}{p}{1.1},
                        just = {\TRule{\False}{\land}[8]}, checked, close
                      ]
                      [\pFmla{\False}{q}{1.1},
                        just = {\TRule{\False}{\land}[8]}, checked
                        [\pFmla{\False}{p}{1.2},
                          just = {\TRule{\False}{\land}[9]}, checked]
                        [\pFmla{\False}{q}{1.2},
                          just = {\TRule{\False}{\land}[9]}, checked, close]
                      ]
                    ]
                  ]
                ]
              ]
            ]
          ]
        ]
      ]
    ]
  \end{oltableau}
  Kari siji cabang bukak, lan cabang iku jangkep. Saka cabang iki
  kita netepake model kanthi himpunan jagad $W = \{1, 1.1, 1.2\}$
  (mung prefiks iki kang katon ing cabang bukak), relasi
  aksesibilitas $R = \{\tuple{1, 1.1}, \tuple{1, 1.2}\}$, lan
  panetepan $V(p) = \{1.1\}$ (amarga larik~6 ngemot
  $\sFmla{\True}{p}[1.1]$) sarta $V(q) = \{1.2\}$ (amarga larik~7
  ngemot $\sFmla{\True}{q}[1.2]$). Model iki kagambar ing
  \olref{fig:counter-Diamond}; bisa dipriksa yen iki kontramodel tumrap
  $(\Diamond p \land \Diamond q) \lif \Diamond (p \land q)$.
  \begin{figure}
  \begin{center}
    \begin{tikzpicture}[modal]
      \node[world] (w1) [label={[align=right]right:\mFalse{p}\\ \mFalse{q}}]{$1$}; 
      \node[world] (w2) [label={[align=right]right:\mTrue{p}\\ \mFalse{q}},
        above left=of w1]{$1.1$}; 
      \node[world] (w3) [label={[align=right]right:\mFalse{p}\\ \mTrue{q}},
        above right=of w1] {$1.2$};
      \draw[->] (w1) to (w2);
      \draw[->] (w1) to (w3);
    \end{tikzpicture}
  \end{center}
  \caption{Kontramodel tumrap $(\Diamond p \land \Diamond q) \lif
    \Diamond (p \land q)$.}
  \ollabel{fig:counter-Diamond}
\end{figure}
\end{ex}          
}

\end{document}
